\documentclass[11pt,a4paper]{article}

% --- Idioma e codificação ---
\usepackage[utf8]{inputenc}
\usepackage[T1]{fontenc}
\usepackage[brazil]{babel}

% --- Layout ---
\usepackage[margin=2.5cm]{geometry}
\usepackage{parskip}

% --- Tabelas, cores e código ---
\usepackage{booktabs}
\usepackage{xcolor}
\usepackage{listings}

% --- Links e metadados ---
\usepackage{hyperref}
\hypersetup{
    colorlinks=true,
    linkcolor=blue!50!black,
    urlcolor=blue!50!black,
    pdftitle={Sucuri --- Guia das coletas manuais pendentes},
    pdfauthor={Projeto Sucuri}
}

\definecolor{codebg}{RGB}{245,245,245}
\lstset{
    backgroundcolor=\color{codebg},
    basicstyle=\ttfamily\small,
    breaklines=true,
    frame=single,
    showstringspaces=false,
    columns=fullflexible,
}

\title{Sucuri --- Guia das coletas manuais pendentes\\
\large O que baixar, onde salvar e o que cada arquivo destrava}
\author{Projeto Sucuri}
\date{21 de julho de 2026}

\begin{document}
\maketitle

\section*{Contexto}

Todo o pipeline automatizável do projeto já roda sozinho: despesas do
Portal da Transparência (com recoleta incremental), IPCA do Banco Central
e população/PIB do IBGE via sidrapy. O que resta são \textbf{fontes sem
API pública utilizável}, que dependem de download ou consulta manual.
Este guia detalha cada uma: passo a passo, nome exato do arquivo de
destino e qual tarefa do \texttt{ROADMAP.md} ela destrava. As mesmas
referências, em forma resumida, estão em \texttt{EXTERNAL.md} (itens
E2--E7).

\begin{table}[h]
\centering
\small
\begin{tabular}{llll}
\toprule
Item & Fonte & Destrava & Esforço estimado \\
\midrule
E3 & INEP (Censo da Educação Superior) & Tarefa 4.1 (custo por aluno) & $\sim$30 min \\
E4 & SIOP (Painel do Orçamento) & Tarefa 4.2 (dotação × execução) & $\sim$30 min \\
E5 & TCU / CGU (consulta caso a caso) & Tarefa 4.3 (validação externa) & horas, iterativo \\
E6 & Portal (download em lote, 1 mês) & Escala das tarefas 3.1/3.7 & $\sim$20 min \\
E7 & API do Portal (reexecução) & Fechar ressalva da tarefa 3.6 & $\sim$30 min de máquina \\
\bottomrule
\end{tabular}
\caption{Visão geral das pendências. Prioridade sugerida: E5 primeiro
(única forma de calibrar a precisão dos sinais já produzidos), depois E3 e E4.}
\end{table}

\section{E3 --- Matrículas por IES (INEP): destrava a tarefa 4.1}

\textbf{Para quê:} calcular \texttt{custo\_por\_aluno} por instituição/ano
(pago real do Conjunto B $\div$ matrículas) e apontar outliers dentro de
cada tipo de instituição. Sem o denominador de matrículas, o ranking de
instituições fica em valor absoluto (ressalva registrada na tarefa 2.1).

\textbf{Passo a passo:}
\begin{enumerate}
    \item Acesse
    \url{https://www.gov.br/inep/pt-br/acesso-a-informacao/dados-abertos/microdados/censo-da-educacao-superior}.
    \item Não é preciso baixar o microdado completo ($\sim$GB): procure a
    \textbf{``Sinopse Estatística da Educação Superior''} do ano desejado
    (arquivo \texttt{.xlsx}, dezenas de MB). A aba de interesse é a de
    matrículas por instituição (IES).
    \item Baixe um arquivo por ano, de 2014 até o último ano publicado
    (o Censo tem defasagem de $\sim$2 anos).
    \item Salve como:
    \texttt{dados/externos/inep/matriculas\_ies\_<ano>.xlsx}
    (criar a pasta \texttt{inep/}; manter um arquivo por ano).
\end{enumerate}

\textbf{Verificação rápida:} a sinopse deve listar as $\sim$110
instituições federais do Conjunto B (universidades e IFs) pelo nome ou
código e-MEC. \textbf{Atenção conhecida:} o nome da IES no INEP não bate
exatamente com o nome do órgão no Portal da Transparência (ex.:
``Universidade Federal Fluminense'' vs.\ ``UFF''); o cruzamento da tarefa
4.1 precisará de uma tabela de correspondência --- o dicionário de siglas
já existente em \texttt{sucuri.utils.sigla\_instituicao} é o ponto de
partida.

\textbf{Depois de salvar:} avisar na sessão de trabalho (ou executar a
tarefa 4.1 do ROADMAP), que produzirá
\texttt{dados/custo\_por\_aluno.parquet} com a ressalva obrigatória de que
o denominador não separa hospital universitário/pesquisa.

\section{E4 --- Dotação orçamentária (SIOP): destrava a tarefa 4.2}

\textbf{Para quê:} comparar a dotação autorizada (LOA + créditos) com o
empenhado do Conjunto A e sinalizar execução $<60\%$ ou $>100\%$ da
dotação por programa/ação --- um sinal de anomalia que a execução sozinha
não revela.

\textbf{Passo a passo:}
\begin{enumerate}
    \item Acesse o Painel do Orçamento:
    \url{https://www1.siop.planejamento.gov.br/paineldoorcamento/}.
    \item No painel principal, aplique os filtros: \textbf{Função = 12
    (Educação)} e \textbf{Subfunção = 364 (Ensino Superior)}; período
    \textbf{2014 até o ano atual}.
    \item Configure a visualização para detalhar por \textbf{Programa} e
    \textbf{Ação} (mesma granularidade do Conjunto A) e inclua as medidas
    \textbf{Dotação Atualizada} (e, se disponível, Dotação Inicial).
    \item Exporte como CSV (botão de exportação do painel).
    \item Salve como: \texttt{dados/externos/siop\_dotacao.csv}.
\end{enumerate}

\textbf{Verificação rápida:} o arquivo deve ter uma linha por
ano×programa×ação com códigos numéricos de programa/ação --- são esses
códigos (não os nomes) que casam com \texttt{codigoPrograma}/
\texttt{codigoAcao} do Conjunto A. Se o painel só exportar nomes,
reexporte incluindo as colunas de código.

\textbf{Depois de salvar:} tarefa 4.2 produz a tabela de divergências
dotação × execução por ano.

\section{E5 --- Acórdãos TCU / relatórios CGU: destrava a tarefa 4.3}

\textbf{Para quê:} a única validação externa possível dos sinais já
produzidos. Sem ela, a precisão dos 176 casos priorizados é desconhecida
por construção. O objetivo é calibrar (``quantos dos top casos aparecem em
achados públicos?''), \textbf{não} ``confirmar culpa''.

\textbf{Passo a passo (iterativo, caso a caso):}
\begin{enumerate}
    \item Abra a lista de casos: \texttt{dados/casos\_priorizados.csv}
    (ou a página ``Anomalias'' do painel Flask, filtrando por 3 sinais).
    \item \textbf{Prioridade máxima:} ``Fundação Euclides da Cunha de
    Apoio Institucional à UFF'' --- único caso atípico em três fontes
    independentes da Fase 3 (concentração de 86,6\% dos contratos da UFF,
    padrão compatível com fracionamento de dispensas e convênios de
    R\$~8,3 milhões). Em seguida: UFRRJ (queda 2019--2021) e UNIVASF
    (salto 2024--2025).
    \item No TCU, pesquise cada instituição em
    \url{https://pesquisa.apps.tcu.gov.br/} (filtrar por acórdãos;
    restringir ao período do sinal ajuda).
    \item Na CGU, pesquise relatórios de auditoria por órgão em
    \url{https://eaud.cgu.gov.br/}.
    \item Para cada caso consultado, anote: instituição, ano do sinal,
    consulta feita, resultado (\textbf{sim} / \textbf{não} /
    \textbf{não pesquisável}) e link/número do acórdão ou relatório.
    \item Salve as notas em:
    \texttt{dados/externos/tcu\_cgu\_notas.md} (formato livre, uma seção
    por caso).
\end{enumerate}

\textbf{Depois de salvar:} tarefa 4.3 consolida a tabela caso × achado
público em \texttt{relatorios/validacao\_tcu.md}. Encontrar ``não'' em
muitos casos é resultado útil (indica falso positivo ou achado ainda não
auditado) --- não força o sinal a ``estar certo''.

\section{E6 --- Mapa código de órgão $\to$ Unidade Gestora: escala das tarefas 3.1 e 3.7}

\textbf{Para quê:} os endpoints \texttt{/despesas/documentos} e
\texttt{/emendas/documentos} identificam a instituição por
\textbf{Unidade Gestora (UG, 6 dígitos SIAFI)}, incompatível com o
\texttt{codigoOrgao} (5 dígitos) usado no resto do projeto, e a API não
tem conversor. Sem esse mapa, a análise por documento só existe para o
piloto (Ouro Preto) e as emendas não são atribuíveis a instituição.

\textbf{Passo a passo:}
\begin{enumerate}
    \item Acesse \url{https://portaldatransparencia.gov.br/download-de-dados}.
    \item Baixe \textbf{um único mês} (qualquer mês recente) do conjunto
    \textbf{``Despesas --- Execução''} (zip com CSVs mensais).
    \item Salve o zip em
    \texttt{dados/externos/portal\_lote/despesas\_execucao/} (manter o
    nome original do arquivo).
    \item O CSV traz \texttt{codigoOrgao} e \texttt{codigoUg} na mesma
    linha: extraia os pares únicos e salve como
    \texttt{dados/externos/mapa\_orgao\_ug.csv} com as colunas
    \texttt{codigoOrgao,codigoUg,nomeOrgao} (a extração pode ser feita na
    sessão de trabalho --- basta o zip estar salvo).
\end{enumerate}

\textbf{Depois de salvar:} \texttt{analises/07\_despesas\_documentos.py}
pode rodar para qualquer instituição do Conjunto B, e as emendas da
tarefa 3.7 passam a ser analisáveis por instituição (hoje: só agregado
nacional/UF).

\section{E7 --- Reexecução do CPGF do EBSERH: fecha a ressalva da tarefa 3.6}

\textbf{Para quê:} a coleta de cartões do EBSERH recebeu HTTP 400 na
página 175 --- erro de servidor que o código trata como fim de paginação;
os 2.610 registros salvos podem estar truncados.

\textbf{Passo a passo (comando longo, rodar fora do horário comercial ---
limite de $\sim$90 req/min):}
\begin{lstlisting}
cd ~/Library/CloudStorage/Dropbox/Repositories/seixas-solutions/sucuri
uv run python - <<'PY'
from sucuri.api import ENV_PATH_PADRAO, carregar_chave_api, criar_sessao
from sucuri.coletores.cartoes import coletar_cartoes
sessao = criar_sessao(carregar_chave_api(ENV_PATH_PADRAO))
# ano a ano para contornar o erro pontual do servidor:
total = 0
for ano in (2023, 2024, 2025):
    registros = coletar_cartoes(sessao, "26443", ano, ano)
    total += len(registros)
    print(ano, len(registros))
print("total EBSERH:", total)
PY
\end{lstlisting}

\textbf{Interpretação:} se o total por janelas anuais superar 2.610, a
coleta original estava truncada --- refazer a análise 3.6 para o EBSERH;
se bater, a ressalva pode ser fechada como ``sem truncamento''.

\section*{Depois de tudo coletado}

Com E3--E5 em \texttt{dados/externos/}, as tarefas 4.1--4.3 do ROADMAP
deixam de estar bloqueadas e podem ser executadas normalmente (cada uma
registra método e resultados em \texttt{relatorios/RELATORIO.md} e na
versão LaTeX, como as demais). A rotina mensal de manutenção do restante
do pipeline segue automatizada --- ver \texttt{EXTERNAL.md}, item X1b.

\textbf{Ressalva permanente:} nada nestas fontes ``confirma''
irregularidade; elas adicionam contexto (custo por aluno, dotação,
achados públicos) para separar atipicidade explicável de indício que
mereça checagem por quem tem competência para tanto.

\end{document}
